\documentclass[10pt,a4paper]{article}
\usepackage[margin=2cm]{geometry}
\usepackage{amsmath,amssymb}
\usepackage{booktabs}
\usepackage{xcolor}
\usepackage{hyperref}
\usepackage[expansion=false]{microtype}
\usepackage[T1]{fontenc}

\hypersetup{colorlinks=true,linkcolor=blue!60!black}

\title{\textbf{The Four-Step Augmentation Pipeline and the\\Exact-Corruption Variant}}
\author{Yani Hammache, Théo Basseras, Titouan Vasnier, Timothy Courtois}
\date{June 2026}

\begin{document}
\maketitle

\section{Where this sits in the pipeline}
\texttt{EEGDataset.\_\_getitem\_\_} (mode \texttt{"ssl"}) reads \emph{one} random
10\,s window $x \in \mathbb{R}^{19 \times 2000}$ (19 channels, 200\,Hz) from a random
recording, z-scores it per channel, then calls \texttt{\_augment(x)} \emph{twice
independently} to produce the two views consumed by the SSL objective:
\[
  v_1 = \texttt{\_augment}(x), \qquad v_2 = \texttt{\_augment}(x).
\]
Neither $v_1$ nor $v_2$ is the clean $x$ — both are corrupted, each with its own draw
of random numbers. The encoder never sees $x$ itself; it only has to learn that $v_1$
and $v_2$ describe ``the same thing'' despite their independent corruptions.

\section{The four steps inside \texttt{\_augment}}
Applied in this fixed order, each gated by its own coefficient:

\begin{enumerate}
  \item \textbf{Amplitude scale jitter.} Per channel $c$, draw
  $s_c \sim \mathcal{U}(1-j,\,1+j)$ with $j = \texttt{aug\_scale\_jitter} = 0.2$, then
  $x_c \leftarrow s_c \cdot x_c$.

  \item \textbf{Additive Gaussian noise.}
  $x \leftarrow x + \varepsilon,\quad \varepsilon \sim \mathcal{N}(0, \sigma^2)$,
  $\sigma = \texttt{aug\_noise\_std} = 0.1$ (z-scored units).

  \item \textbf{Channel dropout.} Each of the 19 channels is independently zeroed
  with probability $p = \texttt{aug\_chan\_drop\_p} = 0.2$:
  $x_c \leftarrow x_c \cdot \mathbb{1}[u_c > p],\ u_c \sim \mathcal{U}(0,1)$.

  \item \textbf{Time masking — legacy vs.\ exact corruption.} Two mutually
  exclusive schemes, selected by \texttt{aug\_exact\_corruption}:

  \begin{description}
    \item[Legacy (\texttt{aug\_exact\_corruption=False}, default).] Zero out
    \emph{one} random contiguous span of length
    $\ell \sim \mathcal{U}(0,\, f)\cdot T$, $f=\texttt{aug\_time\_mask\_frac}=0.2$,
    at a random start position $s \sim \mathcal{U}(0,\, T-\ell)$. At most 20\% of
    the window is touched, in a single block, identical across channels.

    \item[Exact corruption (\texttt{aug\_exact\_corruption=True}).] Per channel,
    independently: pick a random permutation of the $T=2000$ time indices, then
    \begin{itemize}
      \item zero-mask a fixed fraction $m=\texttt{aug\_mask\_frac}=0.2$ of indices,
      \item replace a \emph{disjoint} fraction $o=\texttt{aug\_outlier\_frac}=0.2$
      of indices with outlier spikes $\pm k$ (in z-scored $\sigma$ units,
      $k=\texttt{aug\_outlier\_scale}=6.0$, random sign per spike).
    \end{itemize}
    Exactly $m+o = 40\%$ of each channel's samples are corrupted, scattered
    (not contiguous), and independently chosen per channel.
  \end{description}
\end{enumerate}

\section{What changes between the two schemes}
Steps 1--3 (jitter, noise, channel dropout) are identical in both schemes — only
step 4 differs:

\begin{center}
\begin{tabular}{@{}lll@{}}
\toprule
& \textbf{Legacy} & \textbf{Exact corruption} \\
\midrule
Corrupted fraction & up to 20\%, one block & exactly 40\% \\
Spatial pattern & contiguous span & scattered indices \\
Per-channel independence & same span, all channels & independent per channel \\
Corruption type & zero only & zero \emph{and} $\pm6\sigma$ outliers \\
\bottomrule
\end{tabular}
\end{center}

Both schemes act on the \emph{same} underlying $x$ at the \emph{same} point in the
pipeline (after jitter/noise/dropout, before returning $v_1$/$v_2$) — exact
corruption is not a new data source or an enlarged dataset, it is a stronger,
differently-shaped corruption applied to the one random window already drawn.

\end{document}
